\section{Presenting mathematical concepts inside iPRA}\label{sec:presentations}

\subsection{A common presentation record}

A mathematical concept is embedded by a theorem-specific presentation rather than by assuming a
large ambient universe.

\begin{definition}[Finite presentation interface]
A presentation of a concept \(A\) consists of:
\begin{enumerate}
\item natural-number codes \(a\);
\item a primitive-recursive validity predicate \(\mathsf{Valid}_A(a,\vec q)\);
\item an equality or compatibility relation \(a\sim_A b\), with object proofs of the laws actually
      used;
\item primitive-recursive operations on codes;
\item object derivations of operation preservation and the needed algebraic laws;
\item observation maps to public finite data;
\item producer statements for any code bounds, moduli, covers, or separation facts not derivable by
      the generic verifier.
\end{enumerate}
\end{definition}

The query tuple \(\vec q\) includes every resource that makes the task finite: precision, degree,
cardinality, time horizon, height, or requested observable.  A presentation need not decide equality
of the underlying idealized objects; it needs only the compatibility theorem required by the current
proof.

\subsection{Finite tuples, records, and arrays}

For a fixed number of fields, nested pairing is enough.  The checked pairing module supplies a
primitive-recursive pairing term with proved projections.  A record
\((a_0,\ldots,a_{k-1})\) is represented by a fixed binary tree of pairs.  Every projection theorem is
obtained by composition of \code{fst_pair} and \code{snd_pair}.

For a variable but explicitly bounded horizon \(h\), there are two standard representations.

\paragraph{Mixed-radix code.}
If field \(i\) has a proved bound \(B_i\), encode
\[
  a_0+B_0(a_1+B_1(\cdots)).
\]
Division and remainder recover fields.  The product of the radices is an exact bound for the whole
record.  This is well suited to one-shot bounded search.

\paragraph{Trace code.}
Define the \(i\)-th field by primitive recursion or an array-access function over a code.  A validity
predicate proves the access equations only for \(i<h\).  This representation is better for
incremental construction and avoids enormous explicit radix terms.

A mere cardinality statement is not a code bound.  To search a set of mathematical objects, the
presentation must give a complete encoding whose codes are bounded and whose decoder covers every
valid object.

\subsection{Integers, rationals, and exact algebra}

Integers may be represented by signed pairs or differences of naturals, with equality
\((a,b)\sim(c,d)\) iff \(a+d=b+c\).  Rational codes are pairs \((p,q)\) with \(q>0\); equality is
cross multiplication.  Reduction to lowest terms is optional unless canonical codes are public.

Polynomial codes use a bounded coefficient array and a declared degree.  Evaluation is a bounded
sum; addition and multiplication are bounded convolutions.  Matrices are rectangular arrays, and
matrix multiplication is a bounded sum.  All of these operations are primitive recursive once the
shape parameters are explicit.

The presentation principle is to prove only the quotient laws needed by the theorem.  If a proof
uses rational addition and order but never requests a canonical fraction, there is no need to
formalize normalization.  Equality substitution through observations suffices.

\subsection{Finite sets, maps, and combinatorics}

A finite subset of \([0,N)\) is represented by a bit vector or characteristic term.  Membership,
union, intersection, cardinality, and bounded quantification are primitive recursive.  A finite map
is a table with a validity predicate for key uniqueness; a total map on \([0,N)\) is simply a bounded
array.

Permutations may be represented by a table plus the bounded statement that every value occurs
exactly once.  Graphs are adjacency bits; paths and matchings are finite arrays with local validity
checks.  Counting arguments are expressed by explicit injections/surjections or bounded sums, not by
an abstract powerset.

This representation is directly compatible with the checked pigeonhole, permutation, sums, and
search modules in the derivation library.

\subsection{Algebraic numbers and branches}

An algebraic-number presentation can be a polynomial together with an isolating rational interval
and a certificate that the interval contains exactly one relevant root.  Operations produce new
polynomials and interval certificates.  Equality may be established by a common annihilating
polynomial plus interval overlap/separation, or avoided in favor of stable numerical observations.

A branch label is not constructively available merely because an approximation is accurate.  The
unconditional object is a finite cloud of possible labels or a rational interval.  A public exact
label requires a gap statement showing that distinct labels are separated by more than the current
uncertainty.  In the absence of the gap, \TAPIR{} exports the cloud or invariant observer.

\subsection{Quantitative real and complex names}

A real number used by a fixed theorem is presented by a primitive-recursive rational approximation
\(a(p)\) and an error theorem
\[
  |a(p)-x|\le 2^{-p}
\]
or, without an ambient \(x\), by a Cauchy compatibility theorem
\[
  |a(p)-a(q)|\le 2^{-p}+2^{-q}.
\]
The latter is preferable inside weak arithmetic because it speaks only about finite rational codes.
A complex number is a pair of such names.

The public theorem should be quantitative.  Examples include:
\[
 |F_p(\vec x)-G_p(\vec x)|\le 2^{-p},
\]
``the sign is positive under a supplied gap,'' or ``the output interval has diameter at most
\(2^{-p}\).''  An assertion of exact equality of completed reals is shorthand only when accompanied
by a family of finite error theorems.

\subsection{Continuous functions and moduli}

A continuous function on a coded domain is presented by:

\begin{enumerate}
\item an approximation evaluator \(E(p,x)\);
\item a modulus mapping output precision to input precision;
\item a theorem that compatible input names at the requested input precision produce output
      approximations within the requested error;
\item optional range or Lipschitz bounds.
\end{enumerate}

Without a primitive-recursive modulus, the strongest honest theorem may be finite-consumer or
metastable: for a fixed finite list of input precisions, a supplied certificate meets those requests.
The framework does not turn classical pointwise continuity into a uniform modulus.

Composition is theorem-specialized.  A reverse demand pass asks the outer modulus what precision it
needs from the inner function, constructs the inner approximation at that precision, and then applies
the two error theorems.  The resulting proof graph is acyclic even if the source expression is nested.

\subsection{Differentiation and integration}

A derivative theorem is represented by a finite difference estimate with an explicit remainder:
\[
  \left|\frac{f(x+h)-f(x)}{h}-d(x)\right|\le E(x,h).
\]
A requested precision determines a nonzero rational \(h\) and sufficient input precision.  The
certificate contains the algebraic bound on \(E\).  Exact differentiability is represented by the
family of these finite theorems.

An integral is represented by a quadrature sum plus an error theorem.  Compactness or uniform
continuity enters as a finite mesh/modulus producer.  The checked bounded-sum and Fubini-style modules
handle the finite algebra; the analytic error estimate remains a domain node or residual.

\subsection{Topology and geometry}

For finite topology, use simplicial complexes, nerves, incidence tables, and finite covers.  A
compact-space argument consumes an explicit net or cover rather than invoking unrestricted
compactness.  Homology computations become bounded integer matrix calculations with certificates for
rank, kernels, or Smith data.

Euclidean and projective geometry are represented by coordinates and polynomial incidence
relations.  Degeneracy conditions are explicit.  Noncanonical coordinate charts or orientations are
sealed when the final invariant is chart-independent.  Exact orientation signs require a nonzero
determinant gap; otherwise an interval or disjunction is exported.

\subsection{Probability and finite distributions}

A finite distribution is a list of nonnegative integer weights and a positive total weight.
Expectation and variance are rational bounded sums.  Independence is an equality of finite weight
tables.  Random variables are tables over a finite sample space.

Limit theorems require quantitative tail bounds, rates, or finite-horizon versions.  The finite
algebra can be compiled by existing expectation/variance modules; an asymptotic rate is an explicit
producer.  A statement such as convergence in distribution without a modulus should be kept as a
family of finite test-function obligations or a metastable contract.

\subsection{Groups, categories, and finite diagrams}

A finite group is a multiplication table with identity and inverse witnesses, checked by bounded
universal formulas.  A finitely presented group for a finite-use theorem is represented by words up
to a fixed length plus derivation certificates; unrestricted word-problem claims are not silently
assumed.

A finite category is a table of objects, arrows, source/target, identities, and partial composition,
with bounded associativity checks.  A finite diagram, cone, or natural transformation is another
record.  Universal properties are compiled at the finite diagram actually queried.  Infinite
categories or sheaves enter through finite subdiagrams, covers, and compatibility certificates.

\subsection{Setoids and quotient discipline}

When codes are noncanonical, define a setoid relation \(\sim\) and prove only:

\begin{enumerate}
\item reflexivity, symmetry, and transitivity on valid codes;
\item operation congruence;
\item observer invariance or quantitative compatibility.
\end{enumerate}

The public theorem is stated in invariant observations or in formulas stable under \(\sim\).  There
is no need to construct a quotient type inside first-order arithmetic.  This finite observational
descent is especially effective for bases, charts, presentations, and choices of representatives.

\subsection{A criterion for successful embedding}

A concept is ready for \TAPIR{} compilation when every use in the intended theorem can be reduced to:

\begin{itemize}
\item a fixed finite code or a fixed finite precision request;
\item a primitive-recursive validity/checking predicate;
\item local object derivations for operations and observers;
\item explicit producer boundaries for all non-generic existence or uniformity claims.
\end{itemize}

The criterion is intentionally theorem-relative.  It avoids formalizing a large universe when the
proof consumes only a finite interface.
